\chapter{固定测度的 Token 级优化问题}

\section{外层同步与内层对象}

词表记为 $\mathcal V$，prompt 为 $x$，生成序列为 $y_{1:H}$，第 $t$ 个 token 之前的 prefix 为 $s_t=(x,y_{<t})$。第 $k$ 轮 outer 开始时的冻结 anchor 为 $\pi_k$，本轮第 $n$ 次 inner 更新的 student 为
\begin{equation}
\pi_{k,n}(a\mid s)=\pi_{\theta_{k,n}}(a\mid s),
\qquad \pi_{k,0}=\pi_k.
\end{equation}
教师统一记为 $\pi_T$，在本轮 inner 内冻结。$d_k$ 表示由 $\pi_k$ rollout 诱导的 prefix 访问测度，$\widehat d_k$ 表示当前 rollout batch 的经验测度。inner 固定 prefix batch、anchor 和教师，只更新 student 参数。

\subsection{变长轨迹的访问测度与梯度累积}
\label{sec:prefix-measure-reduction}

固定第 $k$ 轮 rollout 分布。令轨迹 $\tau$ 的有效 response 长度为
$H(\tau)$，包含生成的 EOS，排除 prompt 和 padding；在预算终止时保留实际生成的 token。
假定 $0<\mathbb E H<\infty$。本书明确采用归一化的 token 访问测度
\begin{equation}
 d_k(A)=\frac{\mathbb E_{\tau\sim\pi_k}
 [\sum_{t=1}^{H(\tau)}\mathbf1\{s_t\in A\}]}{\mathbb E H}.
\end{equation}
这里 prompt 也按训练数据分布采样。长轨迹贡献更多被访问的 prefix，
不是先将每条轨迹的总贡献归一化为一。固定 rollout 后，长度和测度均不对 inner 参数求导。

令每个有效位置的损失或 score surrogate 为 $\ell_\theta(s,a)$。
独立 rollout 的经验目标及梯度为
\begin{align}
 N_k&=\sum_{i=1}^{B}H_i,\qquad
 \widehat L_k^{\mathrm{tok}}(\theta)
 =\frac1{N_k}\sum_{i=1}^{B}\sum_{t=1}^{H_i}
 \ell_\theta(s_{i,t},a_{i,t}),\\
 \nabla\widehat L_k^{\mathrm{tok}}
 &=\frac1{N_k}\sum_{i,t}\nabla\ell_\theta(s_{i,t},a_{i,t}).
\end{align}
总体目标为 $\mathbb E_{d_k,\pi_k}\ell_\theta$。有限 batch 的随机分母
使上述比率估计量一般不是总体目标的严格无偏估计；在独立轨迹和可积条件下，
分子、分母的样本均值分别收敛，其比值一致收敛到总体目标。
token 在同一轨迹内相关，因此不能把 $N_k$ 个位置视为独立统计样本来计算置信区间。

\begin{proposition}[访问测度与条件动作分布]
上述 token 访问测度对应的总体 prefix--action 联合分布可分解为
$d_k(s)\pi_k(a\mid s)$。对每条轨迹额外乘以 $1/H(\tau)$ 的序列平均，
一般改变这一条件动作分布。
\end{proposition}
\begin{proof}
在时刻 $t$ 产生动作之前，事件 $\{t\le H\}$ 和 prefix $s_t$
已由历史决定。因此对任意可积函数 $h$，条件期望给出
\begin{align}
 \mathbb E\sum_{t\ge1}\mathbf1\{t\le H\}h(s_t,a_t)
 &=\mathbb E\sum_{t\ge1}\mathbf1\{t\le H\}
 \mathbb E_{a\sim\pi_k(\cdot\mid s_t)}h(s_t,a)\\
 &=\mathbb E H\;\mathbb E_{s\sim d_k,a\sim\pi_k}h(s,a).
\end{align}
除以 $\mathbb E H$ 得到所述分解。预算终止只改变历史可测的存活事件，
不改变该条件期望步骤。

对于序列权重 $1/H$，长度取决于当前及未来生成动作，不能把该权重
从动作条件期望中提出。举例：初始 prefix 下两个动作概率各为 $1/2$；
动作 $a$ 立即 EOS，长度为一；动作 $b$ 后确定生成一个 EOS，长度为二。
按轨迹长度归一化后，初始位置的联合质量分别为 $1/2$ 和 $1/4$，
故给定初始 prefix 的动作概率变成 $2/3$ 与 $1/3$，不同于原策略。
\end{proof}

相比之下，先在每条序列内平均、再在 batch 平均得到
\begin{equation}
 \widehat L_k^{\mathrm{seq}}=
 \frac1B\sum_i\frac1{H_i}\sum_t\ell_\theta(s_{i,t},a_{i,t}).
\end{equation}
同一 token 在两个目标中的权重之比为
$(1/(BH_i))/(1/N_k)=\bar H/H_i$。
例如两条轨迹长度为十和一百，token 平均给两条轨迹的总质量分别为
$1/11$ 与 $10/11$；序列平均给它们各 $1/2$。
这两种目标都可以被定义，但不能在不修改访问测度、条件动作分布和梯度推导的情况下互换。
若每个 prefix 可独立表示、使用完整条件期望且支持集相同，
只改变正的 prefix 权重不改变逐 prefix 的精确 Gibbs 解；
共享参数投影、有限样本和上述依赖未来长度的权重没有这一等价保证。

不除以长度的 batch 序列和 $B^{-1}\sum_i\sum_t\ell$，
则恰等于 $\bar H\widehat L_k^{\mathrm{tok}}$。
在同一冻结 batch 内它是正的常数倍，最优点不变，
但普通梯度尺度、裁剪行为及跨 outer 的相对尺度不同；不能声称 Adam 和裁剪后轨迹也相同。

\begin{proposition}[micro-batch 划分不变性]
将本轮有效位置划分为任意不交 micro-batch $\mathcal B_m$。
若每次 backward 使用
$L_m=N_k^{-1}\sum_{(i,t)\in\mathcal B_m}\ell_{i,t}$，
只在所有 backward 完成后裁剪并执行一次 optimizer step，
则累积梯度与全 batch 的梯度相同，忽略浮点求和次序误差。
\end{proposition}
\begin{proof}
梯度对有限和线性，且冻结分母 $N_k$ 与 inner 参数无关，故
\begin{equation}
 \sum_m\nabla L_m
 =\frac1{N_k}\sum_m\sum_{(i,t)\in\mathcal B_m}\nabla\ell_{i,t}
 =\nabla\widehat L_k^{\mathrm{tok}}.
\end{equation}
若把每个 $L_m$ 改为各自 token 均值后直接累加，
每个位置的权重变成 $1/|\mathcal B_m|$，不再是 $1/N_k$。
只有再乘 $|\mathcal B_m|/N_k$ 才恢复相同梯度。
若每个 micro-batch 单独裁剪，非线性的裁剪映射也不能穿过求和。
\end{proof}

生产实现采用 \texttt{loss\_agg\_mode=token-mean}，
在任何 loss 聚合之前填入当前全 batch 的有效 response token 数。
数据并行若采用梯度平均，每个 rank 的局部和还需乘以数据并行 rank 数，
使跨 rank 平均后恢复全局 token 均值。
每份 micro-batch 不能使用自身长度或上一 outer 的 token 数作为分母。


全篇使用同一组 log-ratio 和 KL 记号：
\begin{align}
U_{k,n}(s,a)&\coloneqq\log\frac{\pi_{k,n}(a\mid s)}{\pi_k(a\mid s)},\\
r_k(s,a)&\coloneqq\log\frac{\pi_T(a\mid s)}{\pi_k(a\mid s)},\\
\mathcal K_k^T(s)&\coloneqq\KL\!\left(\pi_k(\cdot\mid s)\Vert\pi_T(\cdot\mid s)\right),\\
\mathcal K_{k,n}^S(s)&\coloneqq\KL\!\left(\pi_k(\cdot\mid s)\Vert\pi_{k,n}(\cdot\mid s)\right).
\end{align}
因此 $\mathbb E_{\pi_k}[r_k\mid s]=-\mathcal K_k^T(s)$，$\mathbb E_{\pi_k}[U_{k,n}\mid s]=-\mathcal K_{k,n}^S(s)$。前者是冻结 teacher reward 的均值，后者随 student 更新变化。通用 mirror 问题中的信号记为 $f$，其冻结中心化 target 为 $G_k=\eta(f-\mathbb E_{\pi_k}[f\mid s])$。OPD 实例化时取 $f=r_k$，因此
\begin{equation}
G_k(s,a)\coloneqq\eta\bigl(r_k(s,a)+\mathcal K_k^T(s)\bigr),
\qquad\mathbb E_{\pi_k}[G_k\mid s]=0.
\end{equation}
score 向量记为 $\psi_{k,n}(s,a)=\nabla_{\theta_{k,n}}\log\pi_{k,n}(a\mid s)$，其 anchor 条件均值为 $\bar\psi_{k,n}(s)=\mathbb E_{\pi_k}[\psi_{k,n}\mid s]$。

镜像系数为 $\eta>0$，inner 学习率为 $\alpha_{k,n}$。outer 总步数为 $K_{\mathrm{out}}$，每轮 inner 步数为 $N$，rollout batch 大小为 $B$，同一 prefix 的辅助 action 样本数为 $m$，保留的头部 token 数为 $J$。展开截断阶数用 $L$，避免与头部大小混用。所有损失统一最小化，算法使用其梯度 $g_{k,n}=\nabla_{\theta_{k,n}}\mathcal L$ 执行下降。

同步在 inner 结束后进行。hard sync 取 $\pi_{k+1}=\pi_{k,N}$；概率空间的 soft sync 取
\begin{equation}
\pi_{k+1}=(1-\rho)\pi_k+\rho\pi_{k,N},\qquad\rho\in(0,1].
\end{equation}
参数 EMA 是这个概率混合映射的实现近似，需要另行说明。$\lambda_k(s)$ 只表示 PMD-Mean 的归一化乘子，$\gamma$ 表示自然梯度阻尼。student/anchor importance ratio 记为 $\omega_{k,n}=\pi_{k,n}/\pi_k$，K4 的 score surrogate 记为 $c_{k,n}=\pi_{k,n}/\operatorname{sg}(\pi_{k,n})$，其中 $\operatorname{sg}$ 表示停止梯度。

我们研究的 inner 子问题是
\begin{equation}\label{eq:common-objective}
 \max_{p\in\Delta(\mathcal V)}
 \mathcal J_{k,s}(p)
 =\E_{a\sim p}[r_k(s,a)]-
 \frac{1}{\eta}\KL(p\Vert\pi_k),
 \qquad \eta>0.
\end{equation}
这里的 $p$ 是当前 history 上的下一个 token 分布，\emph{不是}整条 response 的分布。访问测度固定只意味着 $s$ 不随 inner 更新；它不意味着优化变量 $p$ 必须等于采样分布 $\pi_k$。

\section{原始 mirror 子问题的 Gibbs 解}

令
\[
 Z_k(s)=\sum_{a\in\mathcal V}\pi_k(a\mid s)\exp(\eta r_k(s,a)).
\]
由于 $r_k=\log\pi_T-\log\pi_k$，有
\begin{equation}\label{eq:token-partition}
 Z_k(s)=\sum_a\pi_k(a\mid s)^{1-\eta}\pi_T(a\mid s)^\eta.
\end{equation}

\begin{theorem}[Token mirror target]
\label{thm:token-target}
若有限词表上 $\pi_k(a\mid s)>0$ 且 $\pi_T(a\mid s)>0$ 对每个 action 成立，则式~\eqref{eq:common-objective} 的唯一最大点为
\begin{equation}\label{eq:token-target}
 \pi_k^{\star,\mathrm{Gibbs}}(a\mid s)=
 \frac{\pi_k(a\mid s)\exp(\eta r_k(s,a))}{Z_k(s)}
 =\frac{\pi_k(a\mid s)^{1-\eta}\pi_T(a\mid s)^\eta}{Z_k(s)}.
\end{equation}
并且对任意 $p$，
\begin{equation}\label{eq:variational-identity}
 \mathcal J_{k,s}(p)
 =\frac{1}{\eta}\log Z_k(s)
 -\frac{1}{\eta}\KL(p\Vert \pi_k^{\star,\mathrm{Gibbs}}).
\end{equation}
\end{theorem}

\begin{proof}
由式~\eqref{eq:token-target}，
\[
 \log \pi_k^{\star,\mathrm{Gibbs}}(a\mid s)=\log\pi_k(a\mid s)+\eta r_k(s,a)-\log Z_k(s).
\]
因此
\begin{align*}
 \KL(p\Vert \pi_k^{\star,\mathrm{Gibbs}})
 &=\sum_a p(a\mid s)\left[
 \log\frac{p(a\mid s)}{\pi_k(a\mid s)}
 -\eta r_k(s,a)+\log Z_k(s)\right]\\
 &=\KL(p\Vert\pi_k)-\eta\E_p[r_k]+\log Z_k(s).
\end{align*}
移项得到式~\eqref{eq:variational-identity}。KL 非负且仅在 $p=\pi_k^{\star,\mathrm{Gibbs}}$ 时为零，所以最大点唯一。\hfill$\square$
\end{proof}

\begin{corollary}[几何插值与极限]
当 $\eta\to0^+$ 时 $\pi_k^{\star,\mathrm{Gibbs}}\to\pi_k$；当 $\eta=1$ 时 $\pi_k^{\star,\mathrm{Gibbs}}=\pi_T$；当 $0<\eta<1$ 时，logits 满足
\[
 \log \pi_k^{\star,\mathrm{Gibbs}}=(1-\eta)\log\pi_k+\eta\log\pi_T-\log Z_k.
\]
因此 $\eta$ 是 token-level teacher interpolation，而不是 inner optimizer 的学习率。
\end{corollary}

\section{Reward 中心化与回归残差}

定义 $\bar r_k(s)=\mathbb E_{a\sim\pi_k}[r_k(s,a)]=-\mathcal K_k^T(s)$。将 reward 换成 $r_k-\bar r_k$ 后，配分函数为
\[
\mathbb E_{\pi_k}[e^{\eta(r_k-\bar r_k)}]=e^{-\eta\bar r_k}Z_k.
\]
分子和配分函数同乘一个 action 无关的常数，原始 mirror 问题的 Gibbs 解不变。下面的平方回归需要区分完整配分函数与均值近似。

三个 residual 全部采用 student 减 target 的方向：
\begin{align}
R^{\mathrm{Part}}_{k,n}(s,a)&\coloneqq U_{k,n}(s,a)-\eta r_k(s,a)+\log Z_k(s),\label{eq:part-residual}\\
R^{\mathrm{Mean}}_{k,n}(s,a)&\coloneqq U_{k,n}(s,a)-G_k(s,a),\label{eq:mean-residual}\\
R^{\mathrm{BPO}}_{k,n}(s,a)&\coloneqq U_{k,n}(s,a)-G_k(s,a)+\mathcal K_{k,n}^S(s).
\end{align}
PMD-Part 的零残差对应 $\pi_{k,n}=\pi_k^{\star,\mathrm{Gibbs}}$。PMD-Mean 以 $\eta\bar r_k$ 近似 $\log Z_k$，仅中心化 reward。对非恒定 reward，由严格 Jensen 不等式有 $\mathbb E_{\pi_k}[e^{G_k}]>1$，所以拟合值 $U=G_k$ 不满足概率归一化，最优残差不能为零。其总体最优分布记为 $\pi_k^{\star,\mathrm{Mean}}$，由第二章的 Lambert-$W$ 解确定。

BPO 则中心化完整的 student--target residual。其 residual 均值为零，而 Mean 的 residual 均值为 $-\mathcal K_{k,n}^S$。这些定义直接区分两个目标，无需额外改变 residual 符号。

\section{参数化策略的投影}

实际神经网络只能表示参数化策略族。固定 prefix 测度 $d_k$ 后，PMD-Part 的总体 inner loss 定义为
\begin{equation}\label{eq:part-loss}
 \mathcal L_{\mathrm{Part}}(\theta_{k,n})
 =\frac12\E_{s\sim d_k,a\sim\pi_k}
 \left[R^{\mathrm{Part}}_{k,n}(s,a)^2\right].
\end{equation}
若模型在每个 history 上可以实现式~\eqref{eq:token-target}，则全局最小值为零；否则最小值是 anchor history 测度下的函数类投影误差。inner steps 增加只能逼近该投影，不能消除模型表达误差或 history coverage 误差。
